%=============================================================================!
% 7.1  FROM VELOCITIES TO MOMENTA                                             !
%=============================================================================!
\section{From velocities to momenta}
\label{sec:ha-legendre}

The state of the bead on its wire (\cref{sec:la-coordinates}) is its
position $x$ and its rate $\dot{x}$, but it can equally be given by $x$
and the momentum $p = \partial\Lag/\partial\dot{x}$ of
\cref{eq:la-momentum}, which for the bead is $m(1 + h'^2)\dot{x}$
(\cref{sec:la-constrained}): since $1 + h'^2$ is positive, each $\dot{x}$ gives
one $p$, and $\dot{x} = p/\bigl(m(1 + h'^2)\bigr)$ gives the rate back.
Hamilton's form of mechanics describes the state by the coordinates and
their momenta. \Cref{sec:ha-liouville} shows that
the motion keeps areas in the plane of $q$ and $p$, which it does not in
general do in the plane of $q$ and $\dot{q}$. The new description needs a
function of $q$ and $p$ that does the work that $\Lag$ does for $q$ and
$\dot{q}$. The tool that exchanges a variable for the slope of a function
of it is the Legendre transform.

\begin{toolbox}[The Legendre transform]
Let $f(v)$ be a function whose slope $f'(v)$ grows steadily with $v$,
that is $f''(v) > 0$ (its graph curves upwards). Then each slope $p$ is
reached at no more than one value of $v$, and we take $p$ among the slopes
that $f$ has. The straight line $pv$ through the origin has that slope
(\cref{fig:ha-legendre}a). The gap between line and
curve, $pv - f(v)$, changes at the rate $p - f'(v)$ as $v$ grows: the
rate is positive while $f'(v) < p$ and negative once $f'(v) > p$, so the
gap rises to its largest value at the $v$ where $f'(v) = p$ and falls
after it. That largest gap is the Legendre transform of $f$,
\begin{equation*}
   f^*(p) = pv - f(v) , \qquad\text{with $v$ such that } f'(v) = p .
\end{equation*}
The tangent to $f$ at that $v$ has slope $p$ and passes through $(v,
f(v))$, so at the vertical axis, a distance $v$ to the left, it is at
$f(v) - pv = -f^*(p)$: the transform is where the tangent of slope $p$
meets the axis, with its sign reversed.

As $p$ changes, $v$ changes with it. By the product rule and the chain
rule,
\begin{align*}
   \frac{\dd f^*}{\dd p}
   &= v + p\,\frac{\dd v}{\dd p} - f'(v)\,\frac{\dd v}{\dd p}
   && \text{differentiating $pv$ and $f(v(p))$}\\
   &= v
   && \text{since $f'(v) = p$ cancels the last two terms.}
\end{align*}
The slope of $f^*$ is $v$, as the slope of $f$ is $p$
(\cref{fig:ha-legendre}b). This slope grows with $p$: differentiating $f'\bigl(v(p)\bigr) = p$ with
respect to $p$ by the chain rule gives $f''(v)\,\dd v/\dd p = 1$, so $\dd
v/\dd p = 1/f''(v) > 0$, as for the inverse functions of
\cref{sec:os-anharmonic}. So $f^*$ can be transformed in its turn. Its
transform at the slope $v$ is $vp - f^*(p)$, with $p$ the slope of $f$ at
which the slope of $f^*$ equals $v$, which is the $p = f'(v)$ we started
from, and
\begin{equation*}
   vp - f^*(p) = vp - \bigl(pv - f(v)\bigr) = f(v) :
\end{equation*}
the transform of the transform is the function we started from, so
$f^*$ carries the same information as $f$, written in terms of the slope.
For example, $f(v) = \frac12 mv^2$ has $f'(v) = mv$, so $p = mv$ and $v =
p/m$, and
\begin{align*}
   f^*(p) &= p\,\frac{p}{m} - \frac12 m\Bigl(\frac{p}{m}\Bigr)^2
   && \text{putting $v = p/m$ into $pv - f(v)$}\\
   &= \frac{p^2}{m} - \frac{p^2}{2m} = \frac{p^2}{2m} ,
\end{align*}
whose slope $p/m$ is indeed $v$.
\end{toolbox}

\begin{figure}[htb]
\centering
\includegraphics{ch07_hamilton/legendre}
\caption{The Legendre transform of $f(v) = \frac12 mv^2$ for a body of
\SI{0.5}{\kilogram}. (a) The line $pv$ with $p =
\SI{1}{\kilogram\metre\per\second}$ (dashed) differs from $f$ by $pv -
f(v)$, which is largest, $f^*(p) = \SI{1}{\joule}$, at $v = p/m =
\SI{2}{\metre\per\second}$, where the slope of $f$ is $p$; the tangent
there (teal) meets the vertical axis at $-f^*(p)$. (b) The transform
$f^*(p) = p^2/(2m)$, whose slope at $p$ is $v$.}
\label{fig:ha-legendre}
\end{figure}

\subsection*{The Hamiltonian}

For one coordinate, the Lagrangian $\Lag(q, \dot{q})$ is a function of
$\dot{q}$ at each fixed $q$, and its slope in $\dot{q}$ is the momentum.
Its Legendre transform in $\dot{q}$, with $q$ held fixed throughout, is
the Hamiltonian,
\begin{equation}
   H(q, p) = p\,\dot{q} - \Lag(q, \dot{q}) ,
   \qquad\text{with $\dot{q}$ such that }
   \frac{\partial\Lag}{\partial\dot{q}} = p ,
   \label{eq:ha-hamiltonian}
\end{equation}
written in terms of $q$ and $p$ alone by putting in $\dot{q}$ found from
$p$; for several coordinates it is $H = \sum_k p_k\dot{q}_k - \Lag$. The
transform needs $\partial\Lag/\partial\dot{q}$ to grow steadily with
$\dot{q}$, which holds for every system met so far, because the kinetic
energy is $K = \frac12 a(q)\,\dot{q}^2$ with a positive $a$: $m$ for a
body on a line, $m(1 + h'^2)$ for the bead, $mL^2$ for the pendulum.

For such a system $p = \partial\Lag/\partial\dot{q} = a\dot{q}$, since $U$
does not contain $\dot{q}$, so $p\dot{q} = a\dot{q}^2 = 2K$, and
\begin{equation*}
   H = 2K - (K - U) = K + U :
\end{equation*}
the Hamiltonian is the energy, now written in terms of the momentum. It
is the quantity $\dot{q}\,\partial\Lag/\partial\dot{q} - \Lag$ that
\cref{ex:la-hamilton} found to be conserved. With several coordinates,
when $K$ is a sum of terms $K_k = \frac12 a_k\dot{q}_k^2$, one for each,
every term gives $p_k\dot{q}_k = a_k\dot{q}_k^2 = 2K_k$, and adding them
again $H = K + U$; \cref{ex:ha-cross}
shows that a kinetic energy with a product of two rates gives the same
result. The examples of Chapter~6 become
\begin{align*}
   &\text{a block on a spring:}
   && p = m\dot{x} ,
   && H = \frac{p^2}{2m} + \frac12 kx^2 ,\\
   &\text{the pendulum:}
   && p = mL^2\dot\theta ,
   && H = \frac{p^2}{2mL^2} + mgL(1 - \cos\theta) ,\\
   &\text{the bead:}
   && p = m(1 + h'^2)\dot{x} ,
   && H = \frac{p^2}{2m\bigl(1 + h'(x)^2\bigr)} + mg\,h(x) ,
\end{align*}
each kinetic energy following from $K = \frac12 a\dot{q}^2$ with $\dot{q}
= p/a$, which gives $K = \frac12 a(p/a)^2 = p^2/(2a)$. The pendulum's
momentum is its angular momentum about the pivot (\cref{sec:ro-plane}).
The bead's kinetic energy now depends on $x$ as well as on $p$, because
the same momentum means a faster bead where the wire is flat than where
it is steep.

For $N$ atoms, with $\vb{p}_i = m_i\dot{\vb{r}}_i$,
\begin{equation*}
   H(\vb{r}^N, \vb{p}^N) = \sum_{i=1}^{N}\frac{\lvert\vb{p}_i\rvert^2}
   {2m_i} + U(\vb{r}^N) ,
\end{equation*}
and the state is the point $\Gamma = (\vb{r}^N, \vb{p}^N)$ of $6N$
numbers. For the puck of \cref{sec:la-polar}, \cref{ex:ha-puck} finds
$H = p_r^2/(2m) + p_\theta^2/(2mr^2) + U(r)$, in which the effective
potential of Chapter~6 appears with $p_\theta$ as the angular momentum $L$
of \cref{sec:la-polar}, not the rod length.

\begin{bridge}
The same transform links the thermodynamic potentials. The free energy of
the electrons at a fixed number $N_{\mathrm{e}}$ has the slope $\mu$, the
chemical potential, with respect to $N_{\mathrm{e}}$; its Legendre
transform in $N_{\mathrm{e}}$, with the sign reversed, is the grand
potential, the free energy minus $\mu N_{\mathrm{e}}$, which is a function
of $\mu$. Grand-canonical DFT holds the electrode's chemical potential
fixed and lets $N_{\mathrm{e}}$ follow, so it minimises the grand
potential, the transform, at fixed $\mu$, as Hamilton's form of mechanics
holds $p$ as the variable in place of $\dot{q}$.
\end{bridge}

\begin{keyidea}
The Hamiltonian $H(q, p) = \sum_k p_k\dot{q}_k - \Lag$ is the Legendre
transform of the Lagrangian in the velocities, written in terms of the
coordinates and the momenta. When the kinetic energy is quadratic in the
velocities it is the energy, $H = K + U$.
\end{keyidea}

\begin{selfcheck}
Write the Hamiltonian of a block of \SI{2}{\kilogram} on a spring of
\SI{8}{\newton\per\metre}. If its energy is \SI{1}{\joule}, what is its
momentum as it passes $x = 0$?
\end{selfcheck}
